\BourbakiExercisesSection

*1. In an affine space \(\mathbf{E}\) over a field \(\mathbf{K}\) , a quadruple \((a, b, c, d)\) of points of \(\mathbf{E}\) is called a parallelogram if \(b - a = c - d\) , in which case \((a, d, c, b)\) is also a parallelogram. Show that if \(\mathbf{K}\) is of characteristic \(\neq 2\) the midpoints of the ordered pairs \((a, c)\) and \((b, d)\) are then equal; what can be said when \(\mathbf{K}\) is of characteristic \(2?_{*}\) .

*2. Let E be an affine space over a field of characteristic \(\neq 2\) and \(a, b, c, d\) any four points of E. Show that if \(x, y, z, t\) are the respective midpoints of the ordered pairs \((a, b), (b, c), (c, d), (d, a)\) , the quadruple \((x, y, z, t)\) is a parallelogram (Exercise 1).*

*3. Let K be a commutative field of characteristic \(\neq 2\) , E an affine plane over K and \(a, b, c, d\) four points of E any three of which do not lie on a straight line. Let \(\mathbf{D}_{xy}\) denote the line passing through two distinct points \(x, y\) of E. Suppose that the lines \(\mathbf{D}_{ab}\) and \(\mathbf{D}_{cd}\) have a common point \(e\) and that \(\mathbf{D}_{ad}\) and \(\mathbf{D}_{bc}\) have a common point \(f\) . Show that the midpoints of the three ordered pairs \((a, c), (b, d), (e, f)\) lie on a straight line. What becomes of this property when \(\mathbf{D}_{ab}\) and \(\mathbf{D}_{cd}\) are parallel or when \(\mathbf{D}_{ad}\) and \(\mathbf{D}_{bc}\) are parallel? Consider the case where K is a field with 3 elements.*

*4. Let K be a field whose characteristic is different from 2 and 3, E an affine space over K, \(a, b, c\) three points of E not on a straight line and \(a', b', c'\) the respective midpoints of the ordered pairs \((b, c)\) , \((c, a)\) and \((a, b)\) . Show that (in the notation of Exercise 3) the lines \(\mathbf{D}_{aa'}\) , \(\mathbf{D}_{bb'}\) and \(\mathbf{D}_{cc'}\) pass through the barycentre of the three points \(a, b, c\) . What becomes of this property when K is of characteristic 2 or of characteristic 3? Generalize to a system of \(n\) affinely independent points.*

\begin{enumerate}
\def\labelenumi{\arabic{enumi}.}
\setcounter{enumi}{4}
\item
  For a non-empty subset V of an affine space E over a field K with at least three elements to be a linear variety, it is necessary and sufficient that for every ordered pair \((x,y)\) of distinct points of V the line \(D_{xy}\) passing through x and y be entirely contained in V. If K has two elements, for V to be a linear variety, it is necessary and sufficient that the barycentre of any three points of V belong to V.
\item
  \begin{enumerate}
  \def\labelenumii{(\alph{enumii})}
  \tightlist
  \item
    Let E be an affine space of dimension \(\geqslant2\) over a field K. For an affine mapping u of E into itself to transform every line of E into a parallel line, it is necessary and sufficient that the linear mapping v associated with u be a homothety \(t\mapsto\gamma t\) of ratio \(\gamma\neq0\) belonging to the centre of K. If \(\gamma=1\) , u is a translation; show that if \(\gamma\neq1\) , there exists one and only one point \(a\in E\) such that \(u(a)=a\) . If a is taken as origin of E, u is then identified with a central homothety for the vector space structure thus determined on E; u is called a central homothety of the affine space E of centre a and ratio \(\gamma\) .
  \end{enumerate}
\end{enumerate}

\begin{enumerate}
\def\labelenumi{(\alph{enumi})}
\setcounter{enumi}{1}
\item
  Let \(u_1, u_2\) be two affine mappings of \(\mathbf{E}\) into \(\mathbf{E}\) each of which is either a translation or a central homothety. Show that \(u_1 \circ u_2\) is a translation or a central homothety of \(\mathbf{E}\) ; if \(u_1, u_2\) and \(u_1 \circ u_2\) are all three central homotheties show that their centres lie on a straight line. What can be said when \(u_1\) and \(u_2\) are central homotheties and \(u_1 \circ u_2\) is a translation?
\item
  Show that the set of translations and central homotheties is a normal subgroup H of the affine group of E and that H/T is isomorphic to the multiplicative group of the centre of K: show that H can only be commutative if H = T, in other words if the centre of K has only two elements.
\end{enumerate}

¶ 7. Let E (resp. E') be an affine space of finite dimension \(n \geqslant 2\) over field K with at least 3 elements (resp. over a field K') and let u be an injective mapping of E into E' mapping any three points in a straight line in E to three points in a straight line and such that the affine linear variety generated by \(u(E)\) in E' is equal to E'.

\begin{enumerate}
\def\labelenumi{(\alph{enumi})}
\item
  Show that \(u\) maps every system of affinely independent points of \(E\) to a system of affinely independent points (use Exercise 5).
\item
  Let \(D_{1}, D_{2}\) be two parallel lines in E and \(D_{1}^{\prime}, D_{2}^{\prime}\) the lines of \(E^{\prime}\) containing respectively \(u(\mathbf{D}_{1})\) and \(u(\mathbf{D}_{2})\) . Show that \(D_{1}^{\prime}\) and \(D_{2}^{\prime}\) are in the same plane; if further u is surjective, show that \(D_{1}^{\prime}\) and \(D_{2}^{\prime}\) are parallel (in the contrary case, show that there would be 3 points not in a straight line in E whose images under u would be in a straight line) (cf.~Exercise 17).
\item
  Suppose henceforth that if \(D_{1}\) and \(D_{2}\) are parallel lines in E, the lines of \(E'\) containing respectively \(u(\mathbf{D}_{1})\) and \(u(\mathbf{D}_{2})\) are parallel. Show that if an origin a in E and the origin \(a' = u(a)\) in \(E'\) are taken, there exists an isomorphism \(\sigma\) of K onto a subfield \(K_{1}\) of \(K'\) such that if E is considered as a vector space over K and \(E'\) as a vector space over \(K_{1}, u\) is an injective semi-linear mapping (relative to \(\sigma\) ) of E into \(E'\) (§ 1, no. 13). (Consider first the case n = 2; given a basis \((e_{1}, e_{2})\) of E, show that for any two elements \(\alpha, \beta\) of K, the points \((\alpha + \beta)e_{1}\) and \((\alpha\beta)e_{1}\) of E can be constructed starting with the points 0, \(e_{1}, e_{2}, \alpha e_{1}, \beta e_{1}\) by constructing parallels to given lines and intersections of given lines; deduce that \(u(\lambda e_{1}) = \lambda^{\sigma}u(e_{1})\) , where \(\sigma\) is an isomorphism of K onto a subfield of \(K'\) , then show that also \(u(\lambda e_{2}) = \lambda^{\sigma}u(e_{2})\) by considering the line joining the points \(\lambda e_{1}\) and \(\lambda e_{2}\) . Pass finally from here to the case where n is arbitrary, arguing by induction on n.) If u is bijective, show that \(K_{1} = K'\) .
\item
  Extend the result of (c) to the case where K is a field with 2 elements assuming also that u maps every system of affinely independent points of E to a system of affinely independent points of E'.
\end{enumerate}

\begin{enumerate}
\def\labelenumi{\arabic{enumi}.}
\setcounter{enumi}{7}
\tightlist
\item
  Let E be a left affine space over a field K and T its translation space.
\end{enumerate}

\begin{enumerate}
\def\labelenumi{(\alph{enumi})}
\tightlist
\item
  For a mapping \(f\) of \(\mathbf{E}\) into a left vector space \(\mathbf{L}\) over \(\mathbf{K}\) to be affine, it is necessary and sufficient that
\end{enumerate}

\[
\begin{array}{r l} f (\mathbf {t} + x) - f (x) & = f (\mathbf {t} + y) - f (y) \\ f (\lambda \mathbf {t} + x) - f (x) & = \lambda f (\mathbf {t} + x) - f (x)) \end{array}
\]

for all \(\lambda\in\mathbf{K},\mathbf{t}\in\mathbf{T},x,y\) in E. Let \(x\mapsto[x]\) be the canonical injection of E into the vector space \(\mathrm{K}_{s}^{(\mathbb{E})}\) of formal linear combinations of elements of E and let N be the subspace of \(\mathrm{K}_{s}^{(\mathbb{E})}\) generated by the elements

\[
[ \mathbf {t} + x ] - [ x ] - [ \mathbf {t} + y ] + [ y ]
\]

\[
[ \lambda \mathbf {t} + x ] - [ x ] - \lambda [ \mathbf {t} + x ] + \lambda [ x ]
\]

for \(\lambda \in \mathbf{K}\) , \(\mathbf{t} \in \mathbf{T}\) , \(x, y\) in \(\mathbf{E}\) ; finally, let \(\mathbf{V}\) be the quotient space of \(\mathbf{K}_s^{(\mathbf{E})}\) by \(\mathbf{N}\) and \(\phi\) the mapping of \(\mathbf{E}\) into \(\mathbf{V}\) which associates with every \(x \in \mathbf{E}\) the class of \([x]\) modulo \(\mathbf{N}\) . Then \(\phi\) is an affine mapping and for every affine mapping \(f\) of \(\mathbf{E}\) into a left vector space \(\mathbf{L}\) over \(\mathbf{K}\) , there exists one and only one linear mapping \(g\) of \(\mathbf{V}\) into \(\mathbf{L}\) such that \(f = g \circ \phi\) .

\begin{enumerate}
\def\labelenumi{(\alph{enumi})}
\setcounter{enumi}{1}
\tightlist
\item
  Let \(\phi_0: \mathbf{T} \to \mathbf{V}\) be the linear mapping associated with \(\phi\) , so that \(\phi(\mathbf{t} + x) - \phi(x) = \phi_0(\mathbf{t})\) . Show that \(\phi\) (and therefore also \(\phi_0\) ) is injective (consider the mapping \(x \mapsto x - a\) of \(\mathbf{E}\) into \(\mathbf{T}\) for some \(a \in \mathbf{E}\) ); for every family \((\lambda_t)_{t \in I}\) of elements of \(\mathbf{K}\) with finite support and every family \((x_t)_{t \in I}\) of elements of \(\mathbf{E}\) ,
\end{enumerate}

\[
\sum_ {i} \lambda_ {i} \phi (x _ {i}) = \phi_ {0} \left(\sum_ {i} \lambda_ {i} x _ {i}\right) \quad \text { if } \quad \sum_ {i} \lambda_ {i} = 0
\]

\[
\sum_ {i} \lambda_ {i} \phi (x _ {i}) = \mu \phi_ {0} \left(\sum_ {i} \mu^ {- 1} \lambda_ {i} x _ {i}\right) \quad \text { if } \quad \sum_ {i} \lambda_ {i} = \mu \neq 0.
\]

Deduce that \(\phi_{0}(T)\) is a hyperplane passing through 0 in V and \(\phi(E)\) a hyperplane parallel to \(\phi_{0}(T)\) .

\begin{enumerate}
\def\labelenumi{\arabic{enumi}.}
\setcounter{enumi}{8}
\tightlist
\item
  Let \(\mathbf{K}\) be a finite field with \(q\) elements and \(\mathbf{V}\) an \(n\) -dimensional vector space over \(\mathbf{K}\) .
\end{enumerate}

\begin{enumerate}
\def\labelenumi{(\alph{enumi})}
\tightlist
\item
  Show that the set of sequences \((x_{1}, x_{2}, \ldots, x_{m})\) of \(m \leqslant n\) vectors of \(\mathbf{V}\) forming a free system has cardinal equal to
\end{enumerate}

\[
(q ^ {n} - 1) (q ^ {n} - q) \dots (q ^ {n} - q ^ {m - 1})
\]

(argue by induction on \(n\) ).

\begin{enumerate}
\def\labelenumi{(\alph{enumi})}
\setcounter{enumi}{1}
\tightlist
\item
  Deduce from (a) that the cardinal of the set of \(m\) -dimensional linear varieties in an \(n\) -dimensional projective space over \(K\) is equal to
\end{enumerate}

\[
\frac {(q ^ {n + 1} - 1) (q ^ {n + 1} - q) \dots (q ^ {n + 1} - q ^ {m})}{(q ^ {m + 1} - 1) (q ^ {m + 1} - q) \dots (q ^ {m + 1} - q ^ {m})}.
\]

\begin{enumerate}
\def\labelenumi{\arabic{enumi}.}
\setcounter{enumi}{9}
\item
  In an \(n\) -dimensional projective space \(\mathbf{P}(\mathbf{V})\) over a field \(\mathbf{K}\) , a projective basis is a set \(\mathbf{S}\) of \(n + 2\) points any two of which form a projectively free system. If \(\mathbf{S} = (a_i)_{0 \leqslant i \leqslant n + 1}\) and \(\mathbf{S}' = (a_i')_{0 \leqslant i \leqslant n + 1}\) are any two projective bases of \(\mathbf{P}(\mathbf{V})\) , show that there exists a transformation \(f \in \mathbf{PGL}(\mathbf{V})\) such that \(f(a_i) = a_i'\) for \(0 \leqslant i \leqslant n + 1\) . For this transformation to be unique, it is necessary and sufficient that \(\mathbf{K}\) be commutative. (Reduce it to the case where \(a_i' = a_i\) for all \(i\) and note that it is always possible to write (in the notation of no. 6) \(a_i = \pi(b_i)\) , where \((b_i)_{1 \leqslant i \leqslant n + 1}\) is a basis of \(\mathbf{V}\) and \(b_0 = b_1 + b_2 + \cdots + b_{n+1}\) .) *Give an example where \(\mathbf{K}\) is a field of quaternions and there exists an infinite set \(\mathbf{T}\) of points of \(\mathbf{P}(\mathbf{V})\) any \(n + 1\) of which form a projectively free system and a transformation \(f \in \mathbf{PGL}(\mathbf{V})\) , distinct from the identity, leaving invariant all the points of \(\mathbf{T}\) .*
\item
  Let V be a 2-dimensional vector space over a field K and a, b, c, d four distinct points of the projective line \(\mathbf{P}(\mathbf{V})\) . The cross-ratio of the quadruple \((a, b, c, d)\) , denoted by \(\begin{bmatrix} a & b \\ d & c \end{bmatrix}\) , is the set of elements \(\xi \in \mathbf{K}\) such that there exist two vectors \(u, v\) in \(\mathbf{V}\) for which (in the notation of no. 6) \(a = \pi(u)\) , \(b = \pi(v)\) , \(c = \pi(u + v)\) , \(d = \pi(u + \xi v)\) . This definition extends immediately to any quadruple of distinct points of a set with a projective line structure (no. 11).
\end{enumerate}

\begin{enumerate}
\def\labelenumi{(\alph{enumi})}
\item
  Show that \(\begin{bmatrix} a & b \\ d & c \end{bmatrix}\) is the set of conjugates of an element \(\neq 1\) of the multiplicative group \(\mathbf{K}^*\) and conversely that, when \(a, b, c\) , are distinct points of \(\mathbf{P}(V)\) and \(\rho\) the set of conjugates of an element \(\neq 1\) of \(\mathbf{K}^*\) , there exists a point \(d \in \mathbf{P}(V)\) such that \(\begin{bmatrix} a & b \\ d & c \end{bmatrix} = \rho\) . For \(d\) to be unique, it is necessary and sufficient that \(\rho\) consist of a single point.
\item
  Show that
\end{enumerate}

\[
\left[ \begin{array}{c c} a & b \\ c & d \end{array} \right] = \left[ \begin{array}{c c} b & a \\ d & c \end{array} \right] = \left[ \begin{array}{c c} a & b \\ d & c \end{array} \right] ^ {- 1}
\]

and

\[
\left[ \begin{array}{c c} d & a \\ c & b \end{array} \right] = 1 - \left[ \begin{array}{c c} a & b \\ d & c \end{array} \right]
\]

(denoting by \(\rho^{-1}\) (resp. \(1 - \rho\) ) the set of conjugates \(\lambda \xi^{-1} \lambda^{-1} = (\lambda \xi \lambda^{-1})^{-1}\) (resp. \(1 - \lambda \xi \lambda^{-1} = \lambda (1 - \xi) \lambda^{-1}\) ), where \(\xi\) is an element of \(\rho\) ).

\begin{enumerate}
\def\labelenumi{(\alph{enumi})}
\setcounter{enumi}{2}
\tightlist
\item
  Let \((a, b, c, d)\) , \((a', b', c', d')\) be two quadruples of distinct points of \(\mathbf{P}(\mathbf{V})\) . For there to exist a bijective semi-linear mapping of V onto itself such that the bijective mapping f of \(\mathbf{P}(\mathbf{V})\) onto itself, obtained when passing to the quotients, satisfies the conditions \(f(a) = a'\) , \(f(b) = b'\) , \(f(c) = c'\) , \(f(d) = d'\) , it is necessary and sufficient that there exist an automorphism \(\sigma\) of K such that
\end{enumerate}

\[
\left[ \begin{array}{c c} a ^ {\prime} & b ^ {\prime} \\ d ^ {\prime} & c ^ {\prime} \end{array} \right] = \left[ \begin{array}{c c} a & b \\ d & c \end{array} \right] ^ {\sigma}.
\]

For there to exist a transformation f of the projective group PGL(V) satisfying the above conditions, it is necessary and sufficient that

\[
\left[ \begin{array}{c c} a ^ {\prime} & b ^ {\prime} \\ d ^ {\prime} & c ^ {\prime} \end{array} \right] = \left[ \begin{array}{c c} a & b \\ d & c \end{array} \right].
\]

\begin{enumerate}
\def\labelenumi{\arabic{enumi}.}
\setcounter{enumi}{11}
\tightlist
\item
  Let \(\mathbf{P}(\mathbf{V})\) be a (left) projective space of finite dimension \(n\) over a field \(\mathbf{K}\) . Show that there exists on the set of projective hyperplanes of \(\mathbf{P}(\mathbf{V})\) an \(n\) -dimensional (right) projective space structure over \(\mathbf{K}\) canonically isomorphic to that of \(\mathbf{P}(\mathbf{V}^*)\) ( \(\mathbf{V}^*\) being the dual of \(\mathbf{V}\) ). If \(\mathbf{M}\) is a linear variety of dimension \(r < n\) in \(\mathbf{P}(\mathbf{V})\) , derive an \((n - r - 1)\) -dimensional projective space structure on the set of projective hyperplanes containing \(\mathbf{M}\) . In particular if \(\mathbf{M}\) is of dimension \(n - 2\) , the cross-ratio \(\begin{bmatrix} \mathrm{H}_1 & \mathrm{H}_2 \\ \mathrm{H}_4 & \mathrm{H}_3 \end{bmatrix}\) of a quadruple \((\mathrm{H}_1, \mathrm{H}_2, \mathrm{H}_3, \mathrm{H}_4)\) of distinct hyperplanes containing M can be defined. Show that if \(\mathbf{D} \subset \mathbf{P}(\mathbf{V})\) is a line not meeting M and \(a_i\) is the intersection of D with \(\mathrm{H}_i\) ( \(1 \leqslant i \leqslant 4\) ), then
\end{enumerate}

\[
\left[ \begin{array}{c c} a _ {1} & a _ {2} \\ a _ {4} & a _ {3} \end{array} \right] = \left[ \begin{array}{c c} \mathrm{H} _ {1} & \mathrm{H} _ {2} \\ \mathrm{H} _ {4} & \mathrm{H} _ {3} \end{array} \right].
\]

*13. In a projective plane \(\mathbf{P}(\mathbf{V})\) over a field \(\mathbf{K}\) of characteristic \(\neq 2\) , let \(a, b, c, d\) be four points forming a projective basis (Exercise 10); denoting by \(\mathbf{D}_{xy}\) the line passing through two distinct points \(x, y\) of \(\mathbf{P}(\mathbf{V})\) , let \(e, f, g\) be the points of intersection of the lines \(\mathbf{D}_{ab}\) and \(\mathbf{D}_{cd}\) , \(\mathbf{D}_{ac}\) and \(\mathbf{D}_{bd}\) , \(\mathbf{D}_{ad}\) and \(\mathbf{D}_{bc}\) respectively; let \(h\) be the point of intersection of \(\mathbf{D}_{bc}\) and \(\mathbf{D}_{ef}\) ; show that \(\begin{bmatrix} b & c \\ h & g \end{bmatrix} = \{-1\}\) (``theorem of the complete quadrilateral''; reduce it to the case where \(\mathbf{D}_{ad}\) is the line at infinity of an affine plane). What is the corresponding result when \(\mathbf{K}\) is of characteristic \(2?_{*}\)

\begin{enumerate}
\def\labelenumi{\arabic{enumi}.}
\setcounter{enumi}{13}
\item
  In a projective plane \(\mathbf{P}(\mathrm{V})\) over a field \(\mathbf{K}\) with at least three elements, let \(\mathbf{D}, \mathbf{D}'\) be two distinct lines. In order that, for any distinct points \(a, b, c\) of \(\mathbf{D}, a', b', c'\) of \(\mathbf{D}'\) , the intersection points \(r\) of \(\mathbf{D}_{ab'}\) and \(\mathbf{D}_{ba'}\) , \(q\) of \(\mathbf{D}_{ac'}\) and \(\mathbf{D}_{ca'}\) , \(p\) of \(\mathbf{D}_{bc'}\) and \(\mathbf{D}_{cb'}\) lie on a straight line, it is necessary and sufficient that \(\mathbf{K}\) be commutative (``Pappus's theorem''; reduce it to the case where \(q\) and \(r\) are on the line at infinity of an affine plane). Apply this theorem to the projective space of lines of \(\mathbf{P}(\mathrm{V})\) (Exercise 12).
\item
  In a projective plane over a field K with at least three elements, let \(s, a, b, c, a', b', c'\) be seven distinct points such that \(\{s, a, b, c\}\) and \(\{s, a', b', c'\}\) are projective bases (Exercise 10) and the lines \(\mathbf{D}_{sa}, \mathbf{D}_{sb}, \mathbf{D}_{sc}\) pass respectively through \(a', b', c'\) . Show that the intersection points \(r\) of \(\mathbf{D}_{ab}\) and \(\mathbf{D}_{a'b'}\) , \(p\) of \(\mathbf{D}_{bc}\) and \(\mathbf{D}_{b'c'}\) , \(q\) of \(\mathbf{D}_{ca}\) and \(\mathbf{D}_{c'a'}\) lie on a straight line (``Desargues's theorem''; method analogous to that of Exercise 14).
\item
  \begin{enumerate}
  \def\labelenumii{(\alph{enumii})}
  \tightlist
  \item
    Let \(\mathbf{E} = \mathbf{P}(\mathbf{V})\) and \(\mathbf{E}' = \mathbf{P}(\mathbf{V}')\) be two projective spaces of the same dimension \(n \geqslant 2\) over two fields \(\mathbf{K}\) , \(\mathbf{K}'\) respectively and let \(u\) be a bijective mapping of \(\mathbf{E}\) onto \(\mathbf{E}'\) , mapping any three points on a straight line to three points on a straight line. Show that there exist an isomorphism \(\sigma\) of \(\mathbf{K}\) onto \(\mathbf{K}'\) and a bijective semi-linear mapping \(v\) of \(\mathbf{V}\) onto \(\mathbf{V}'\) (relative to \(\sigma\) ) such that \(u\) is the mapping obtained from \(v\) when passing to the quotients (``fundamental theorem of projective geometry''; use Exercise 7). Suppose also that \(\mathbf{V}' = \mathbf{V}\) and \(\mathbf{K}\) is commutative; for \(u\) to be a projective mapping, it is necessary and sufficient that also \(\begin{bmatrix} u(a) & u(b) \\ u(d) & u(c) \end{bmatrix} = \begin{bmatrix} a & b \\ d & c \end{bmatrix}\) for every quadruple \((a, b, c, d)\) of distinct points on a straight line in \(\mathbf{P}(\mathbf{V})\) .
  \end{enumerate}
\end{enumerate}

\begin{enumerate}
\def\labelenumi{(\alph{enumi})}
\setcounter{enumi}{1}
\tightlist
\item
  Let \(p\) be an integer such that \(1 \leqslant p \leqslant n - 1\) . Show that the first conclusion of (a) holds when the image under \(u\) of every \(p\) -dimensional projective linear variety is contained in a \(p\) -dimensional projective linear variety.
\end{enumerate}

\begin{enumerate}
\def\labelenumi{\arabic{enumi}.}
\setcounter{enumi}{16}
\tightlist
\item
  Let V be a vector space of finite dimension n over a field K, \((e_{i})_{1\leq i\leq n}\) a basis of V, \(K'\) a subfield of K and \(V'\) the n-dimensional vector space over \(K'\) generated by the \(e_{i}\) . Give an example of an injective mapping of \(V'\) into V, mapping any three points on a straight line of the affine space \(V'\) to three points on a straight line in the affine space V, but not necessarily mapping two parallel lines to sets contained in two parallel lines. (Embed V in the projective space E canonically associated with it and consider a projective transformation u of E into itself such that the inverse image under u of the hyperplane at infinity is distinct from this hyperplane and contains no point of \(V'\) ; for example K can be taken to be infinite and \(K'\) finite.)
\end{enumerate}

¶ 18. Let \(\mathbf{E} = \mathbf{P}(\mathbf{V})\) be a projective plane over a field \(\mathbf{K}\) and \(u\) a bijective mapping of \(\mathbf{E}\) onto itself, arising when passing to the quotients from a bijective semi-linear mapping \(v\) of \(\mathbf{V}\) onto itself, relative to an automorphism \(\sigma\) of \(\mathbf{K}\) .

\begin{enumerate}
\def\labelenumi{(\alph{enumi})}
\item
  Show that the four following properties are equivalent: (α) for all \(x \in E\) , \(x, u(x)\) and \(u^{2}(x)\) lie on a straight line; (β) every line of E contains a point invariant under u; (γ) through every point of E there passes a line invariant under u; (δ) for every line D of E, the three line D, \(u(\mathbf{D})\) and \(u^{2}(\mathbf{D})\) have a common point. (Show first that (α) and (γ) are equivalent; deduce by duality (Exercise 12) that (β) and (δ) are equivalent; prove finally that (γ) implies (β) and deduce by duality that (β) implies (γ).)
\item
  Suppose that u has the properties stated in (a). Show that if there exists in E a line D invariant under u and containing only a single point invariant under u, u arises from a transvection v of V when passing to the quotients (§ 10, Exercise 11). (Show that every line invariant under u contains a; by considering a line not passing through a, show that there exists a line \(D_{0}\) passing through a and containing at least two points invariant under u; conclude that all the points of \(D_{0}\) are necessarily invariant under u.)
\item
  Suppose that u has the properties stated in (a). Show that if there exists in E a line E invariant under u and containing only two points invariant under u, u arises from a dilatation v of V when passing to the quotients ( \(\S\) 10, Exercise 11). (If a, b are the two points of D invariant under u, show that every line invariant under u passes through a or through b; then note that there exist at least two other points c, d distinct from a, b and invariant under u and therefore the line \(D_{cd}\) passes through a or b; conclude by proving that all the points of \(D_{cd}\) are invariant under u.)
\item
  Suppose that u has the properties stated in (a) and that every line of E invariant under u contains at least three distinct points invariant under u; then there exists in E a projective basis (Exercise 10) each point of which is invariant under u; conclude that there exists a basis \((e_{i})_{1\leq i\leq3}\) of V such that u arises from a semi-linear mapping \(v\) of \(\mathbf{V}\) into itself such that \(v(e_i) = e_i\) for \(1 \leqslant i \leqslant 3\) , when passing to the quotients. The set of points of \(\mathbf{E}\) invariant under \(u\) is then the projective plane \(\mathbf{P}(\mathbf{V}_0)\) , where \(\mathbf{V}_0\) is the vector space over the field \(\mathbf{K}_0\) of invariants of \(\sigma\) , generated by \(e_1, e_2, e_3\) .
\item
  Suppose henceforth that \(u\) satisfies the conditions of (a) and (d) and that neither \(u\) nor \(u^2\) is the identity. Show that there exists \(\gamma \in \mathbf{K}\) such that \(\gamma^\sigma = \gamma\) and
\end{enumerate}

\[
(\xi^ {\sigma} - \xi) ^ {\sigma} = \gamma (\xi^ {\sigma} - \xi)\tag{1}
\]

for all \(\xi \in \mathbf{K}\) (use conditions \((\alpha)\) of \((a)\) and the existence of \(\zeta \in \mathbf{K}\) such that \(\zeta^{\sigma} \neq \zeta\) ). Show that \(\gamma \neq -1\) and that, for all \(\xi \in \mathbf{K}\) such that \(\xi^{\sigma} \neq \xi\) ,

\[
(1 + \gamma) \xi^ {\sigma} \gamma = \gamma \xi (1 + \gamma)\tag{2}
\]

(apply (1) replacing \(\xi\) by \(\xi^2\) ); extend (2) to all \(\xi \in \mathbf{K}\) by noting that \(\xi = \eta - \zeta\) , where \(\eta^\sigma \neq \eta\) and \(\xi^\sigma \neq \zeta\) . Conclude that \(\gamma \neq 1\) , then deduce from (1) and (2) that

\[
\xi^ {\sigma} = (1 + \gamma) \xi (1 + \gamma) ^ {- 1}\tag{3}
\]

for all \(\xi \in K\) and that \(\gamma + \gamma^{-1}\) belongs to the centre of \(K\) . Obtain the converse. *Give an example where \(\gamma\) does not belong to the centre of \(K\) , but where \(\gamma + \gamma^{-1}\) is in this centre.*

*19. Let K be a commutative field, \(\tilde{K}\) the projective field obtained by adjoining to K a point at infinity (no. 9) and \(f\) and \(g\) two elements of K(X). Show that, writing \(h(\mathbf{X}) = f(g(\mathbf{X}))\) , if \(\tilde{f}, \tilde{g}, \tilde{h}\) are the canonical extensions of \(f, g, h\) to \(\tilde{\mathbf{K}}\) , then \(\tilde{h} = \tilde{f} \circ \tilde{g}_{\cdot *}\)

